% ============================================================
\chapter{Espaces $T_0$ et Ordre de Sp\'ecialisation}
\label{ch:t0_ordre}
% ============================================================

L'une des d\'ecouvertes les plus f\'econdes de la topologie du XX\ieme{}
si\`ecle est que derri\`ere chaque espace topologique se cache un
\emph{ordre}. Cet ordre, dit \og de sp\'ecialisation\fg, relie la topologie
\`a la th\'eorie des ordres partiels de mani\`ere profonde et inattendue.
Dans le monde des espaces $T_0$, cette correspondance devient une
\'equivalence~: la topologie \emph{est} l'ordre, et l'ordre \emph{est} la
topologie. C'est cette dualit\'e qui fait des espaces $T_0$ les objets
centraux de ce cours.

\begin{intuition}
L'ordre de sp\'ecialisation est le pont fondamental entre topologie et th\'eorie de l'ordre.
Dans un espace $T_0$, chaque point poss\`ede une <<~position~>> dans un ordre partiel
d\'etermin\'e par la topologie. Ce chapitre explore syst\'ematiquement cette correspondance
et ses cons\'equences.
\end{intuition}

% ------------------------------------------------------------
\section{L'ordre de sp\'ecialisation}
% ------------------------------------------------------------

\begin{definition}[Ordre de sp\'ecialisation]
Soit $(X,\tau)$ un espace topologique. L'\emph{ordre de sp\'ecialisation} est la relation
$\leq$ sur $X$ d\'efinie par~:
\[
  x \leq y \quad \Longleftrightarrow \quad x \in \overline{\{y\}}
  \quad \Longleftrightarrow \quad \text{tout ouvert contenant } x \text{ contient aussi } y.
\]
\end{definition}

\begin{proposition}[Propri\'et\'es de base]
Soit $(X,\tau)$ un espace topologique et $\leq$ son ordre de sp\'ecialisation.
\begin{enumerate}
  \item $\leq$ est un pr\'eordre (r\'eflexif et transitif).
  \item $\leq$ est un ordre partiel (antisym\'etrique) si et seulement si $X$ est $T_0$.
  \item $x \leq y$ si et seulement si $\overline{\{y\}} \subseteq \overline{\{x\}}$
    (attention \`a l'inversion !).
  \item Tout ouvert $U$ est une partie haute~: si $x \in U$ et $x \leq y$, alors $y \in U$.
  \item Tout ferm\'e $F$ est une partie basse~: si $y \in F$ et $x \leq y$, alors $x \in F$.
\end{enumerate}
\end{proposition}

\begin{proof}
(1) R\'eflexivit\'e~: $x \in \overline{\{x\}}$ toujours. Transitivit\'e~: si $x \leq y$ et
$y \leq z$, alors pour tout ouvert $U \ni x$, on a $y \in U$ (car $x \leq y$), puis
$z \in U$ (car $y \leq z$), donc $x \leq z$.

(2) Antisym\'etrie~: $x \leq y$ et $y \leq x$ signifient $\overline{\{x\}} =
\overline{\{y\}}$, c'est-\`a-dire que $x$ et $y$ sont topologiquement indistinguables.
C'est \'equivalent \`a $x = y$ ssi $X$ est $T_0$.

(3) $x \leq y$ signifie $x \in \overline{\{y\}}$, donc $\overline{\{x\}} \subseteq
\overline{\{y\}}$... non, attention~! $x \leq y$ signifie que $y$ sp\'ecialise $x$, i.e.
$x \in \overline{\{y\}}$. L'adh\'erence de $\{y\}$ est $\downarrow y$ (les points
sp\'ecialis\'es par $y$).

(4) Soit $U$ ouvert, $x \in U$, $x \leq y$. Tout ouvert contenant $x$ contient $y$, en
particulier $U$, donc $y \in U$.

(5) Analogue par passage au compl\'ementaire.
\end{proof}

\begin{attention}
La convention d'ordre varie selon les auteurs. Nous suivons la convention o\`u
$x \leq y$ signifie <<~$x$ se sp\'ecialise en $y$~>>, c'est-\`a-dire $x \in \overline{\{y\}}$,
ou encore $y$ est <<~plus g\'en\'erique~>> que $x$. Certains auteurs utilisent la convention
oppos\'ee. Dans notre convention, les ouverts sont des parties \emph{hautes}.
\end{attention}

% ------------------------------------------------------------
\section{Topologie d'Alexandrov}
% ------------------------------------------------------------

\begin{definition}[Topologie d'Alexandrov]
Soit $(P, \leq)$ un ensemble pr\'eordonn\'e. La \emph{topologie d'Alexandrov} sur $P$ est
la topologie dont les ouverts sont exactement les parties hautes (\emph{upward-closed})~:
\[
  U \in \tau_{\mathrm{Alex}} \quad \Longleftrightarrow \quad
  (x \in U \text{ et } x \leq y \Rightarrow y \in U).
\]
\end{definition}

\begin{theoreme}[Caract\'erisation de la topologie d'Alexandrov]
Un espace topologique $(X,\tau)$ porte la topologie d'Alexandrov de son ordre de
sp\'ecialisation si et seulement si $\tau$ est stable par intersections quelconques.
\end{theoreme}

\begin{proof}
$(\Rightarrow)$ Les parties hautes d'un pr\'eordre sont stables par intersections
quelconques.

$(\Leftarrow)$ Si $\tau$ est stable par intersections quelconques, alors pour tout $x$,
l'intersection $\bigcap\{U \in \tau : x \in U\}$ est un ouvert, not\'e $U_x$. C'est le
plus petit ouvert contenant $x$. On a $y \in U_x$ si et seulement si tout ouvert contenant
$x$ contient $y$, c'est-\`a-dire $x \leq y$. Ainsi $U_x = \uparrow x$ et tout ouvert est
une r\'eunion de tels ensembles, donc une partie haute.
\end{proof}

\begin{proposition}[\'Equivalence de cat\'egories]
Le foncteur qui associe \`a un ensemble pr\'eordonn\'e $(P,\leq)$ l'espace d'Alexandrov
$(P, \tau_{\mathrm{Alex}})$ \'etablit une \'equivalence de cat\'egories entre
$\mathbf{Preord}$ (pr\'eordres et fonctions croissantes) et la sous-cat\'egorie pleine
de $\mathbf{Top}$ form\'ee des espaces d'Alexandrov. La restriction aux ordres partiels
donne une \'equivalence avec les espaces d'Alexandrov $T_0$.
\end{proposition}

\begin{proof}
Une application $f\colon P \to Q$ est croissante si et seulement si la pr\'eimage de
toute partie haute est haute, ce qui \'equivaut \`a la continuit\'e pour les topologies
d'Alexandrov. Le foncteur inverse associe \`a un espace d'Alexandrov son ordre de
sp\'ecialisation.
\end{proof}

% ------------------------------------------------------------
\section{Adh\'erence et ordre}
% ------------------------------------------------------------

\begin{proposition}[Adh\'erence et parties basses]
Dans tout espace topologique $(X,\tau)$~:
\begin{enumerate}
  \item $\overline{\{x\}} = \downarrow\! x$ (pour l'ordre de sp\'ecialisation).
  \item Pour tout $A \subseteq X$, $\overline{A} = \downarrow\! A$ si et seulement si
    $X$ porte la topologie d'Alexandrov.
  \item En g\'en\'eral, $\downarrow\! A \subseteq \overline{A}$, avec \'egalit\'e pour
    les singletons.
\end{enumerate}
\end{proposition}

\begin{proof}
(1) Par d\'efinition, $y \in \overline{\{x\}}$ si et seulement si tout ouvert contenant $y$
contient $x$, c'est-\`a-dire $y \leq x$, soit $y \in \downarrow\! x$.

(3) Si $a \in A$ et $y \leq a$, alors $y \in \overline{\{a\}} \subseteq \overline{A}$,
donc $\downarrow\! A \subseteq \overline{A}$.
\end{proof}

\begin{definition}[Point g\'en\'erique]
Soit $F$ un ferm\'e irr\'eductible de $X$ (c'est-\`a-dire $F \neq \emptyset$ et si
$F = F_1 \cup F_2$ avec $F_1, F_2$ ferm\'es, alors $F = F_1$ ou $F = F_2$). Un point
$x \in F$ est un \emph{point g\'en\'erique} de $F$ si $\overline{\{x\}} = F$.
\end{definition}

\begin{remarque}
L'existence et l'unicit\'e de points g\'en\'eriques pour les ferm\'es irr\'eductibles sont
au c\oe ur de la notion d'espace sobre (Chapitre~3).
\end{remarque}

% ------------------------------------------------------------
\section{Continuit\'e et monotonie}
% ------------------------------------------------------------

\begin{theoreme}[Continuit\'e implique monotonie]
Si $f\colon X \to Y$ est continue, alors $f$ est croissante pour les ordres de
sp\'ecialisation~: $x \leq_X y \Rightarrow f(x) \leq_Y f(y)$.
\end{theoreme}

\begin{proof}
Soit $V$ un ouvert de $Y$ contenant $f(x)$. Alors $f^{-1}(V)$ est un ouvert de $X$
contenant $x$. Comme $x \leq y$, on a $y \in f^{-1}(V)$, donc $f(y) \in V$. Ainsi
$f(x) \leq_Y f(y)$.
\end{proof}

\begin{attention}
La r\'eciproque est fausse en g\'en\'eral~: une fonction croissante entre espaces $T_0$
n'est pas n\'ecessairement continue. La continuit\'e est une condition \emph{plus forte}
que la monotonie. Cependant, pour les topologies d'Alexandrov, les deux notions
co\"incident.
\end{attention}

\begin{exemple}
Soit $X = \{0,1\}$ avec la topologie discr\`ete et $Y = \mathbb{S}$ l'espace de
Sierpi\'nski. L'ordre de sp\'ecialisation de $X$ est l'\'egalit\'e, celui de $Y$ est
$1 \leq 0$. La fonction $f(0) = 1$, $f(1) = 0$ est croissante (trivialement, car
l'ordre sur $X$ est discret) mais elle est continue. La fonction constante $g \equiv 0$
est \'egalement croissante et continue. Par contre, consid\'erons $X = \{a,b\}$ avec
$\tau = \{\emptyset, \{a\}, \{a,b\}\}$ (donc $b \leq a$) et $Y = \{c,d\}$ avec
$\tau' = \{\emptyset, \{c\}, \{c,d\}\}$ (donc $d \leq c$). La fonction $f(a)=d$,
$f(b)=c$ est croissante ($b \leq a$ donne $f(b)=c \geq d = f(a)$... non, $d \leq c$
donc $f(a) \leq f(b)$), et elle n'est pas continue car $f^{-1}(\{c\}) = \{b\}$ n'est
pas ouvert.
\end{exemple}

% ------------------------------------------------------------
\section{Le foncteur de sp\'ecialisation}
% ------------------------------------------------------------

\begin{definition}[Foncteur de sp\'ecialisation]
Le \emph{foncteur de sp\'ecialisation} $\Sigma\colon \mathbf{Top}_0 \to \mathbf{Pos}$
associe \`a chaque espace $T_0$ son ensemble sous-jacent muni de l'ordre de
sp\'ecialisation, et \`a chaque application continue la m\^eme fonction (qui est croissante
par le th\'eor\`eme pr\'ec\'edent).
\end{definition}

\begin{theoreme}[Adjoint \`a gauche du foncteur de sp\'ecialisation]
Le foncteur d'Alexandrov $\mathrm{Alex}\colon \mathbf{Pos} \to \mathbf{Top}_0$ est adjoint
\`a gauche du foncteur de sp\'ecialisation $\Sigma$~:
\[
  \mathrm{Alex} \dashv \Sigma.
\]
Autrement dit, pour tout poset $P$ et tout espace $T_0$ $X$~:
\[
  \mathrm{Hom}_{\mathbf{Top}_0}(\mathrm{Alex}(P), X) \cong
  \mathrm{Hom}_{\mathbf{Pos}}(P, \Sigma(X)).
\]
\end{theoreme}

\begin{proof}
Un morphisme $\mathrm{Alex}(P) \to X$ dans $\mathbf{Top}_0$ est une application continue,
donc en particulier croissante de $(P,\leq)$ vers $(\Sigma(X), \leq_X)$. R\'eciproquement,
une application croissante $f\colon P \to \Sigma(X)$ est continue $\mathrm{Alex}(P) \to X$
car la pr\'eimage d'un ouvert $U$ de $X$ (partie haute pour $\leq_X$) est une partie haute
de $P$ (car $f$ est croissante), donc un ouvert d'Alexandrov.
\end{proof}

% ------------------------------------------------------------
\section{Espaces $T_0$ et $d$-espaces}
% ------------------------------------------------------------

\begin{definition}[$d$-espace]
Un espace topologique $X$ est un \emph{$d$-espace} (ou espace monotoniquement convergent)
si pour tout ensemble dirig\'e $D$ dans $(X, \leq)$ (pour l'ordre de sp\'ecialisation)
ayant un supremum $\bigvee D$ dans $(X,\leq)$, les ouverts contenant $\bigvee D$
contiennent un \'el\'ement de $D$.
\'Equivalemment, si $D$ est un sous-ensemble dirig\'e et $\overline{D}$ est irr\'eductible,
alors $\overline{D}$ a un point g\'en\'erique qui est le supremum de $D$.
\end{definition}

\begin{theoreme}
Tout espace sobre est un $d$-espace. Tout espace $T_1$ est un $d$-espace (trivialement,
car les ensembles dirig\'es pour l'ordre discret sont les singletons).
\end{theoreme}

\begin{proof}
Soit $X$ sobre et $D$ un ensemble dirig\'e ayant un supremum $s = \bigvee D$.
L'adh\'erence $\overline{D}$ est un ferm\'e, et $\overline{D} = \downarrow\!\overline{D}$
(tout ferm\'e est une partie basse). Montrons que $\overline{D}$ est irr\'eductible.
Si $\overline{D} = F_1 \cup F_2$ avec $F_1, F_2$ ferm\'es, alors $D \subseteq F_1 \cup F_2$.
Comme $D$ est dirig\'e et les $F_i$ sont ferm\'es (donc parties basses), si $D \not\subseteq
F_1$ et $D \not\subseteq F_2$, il existe $d_1 \in D \setminus F_1$ et $d_2 \in D \setminus
F_2$. Par directivit\'e, il existe $d \geq d_1, d_2$ dans $D$. Alors $d \in F_1 \cup F_2$,
disons $d \in F_1$. Comme $F_1$ est une partie basse et $d_1 \leq d$, on a $d_1 \in F_1$,
contradiction. Donc $\overline{D}$ est irr\'eductible. Par sobri\'et\'e, $\overline{D}$
a un point g\'en\'erique, qui est n\'ecessairement $s$.
\end{proof}

% ------------------------------------------------------------
\section{Espaces d'Alexandrov et treillis finis}
% ------------------------------------------------------------

\begin{exemple}[Topologie d'Alexandrov sur un treillis fini]
Soit $L$ un treillis fini. La topologie d'Alexandrov a pour ouverts les parties hautes.
L'ordre de sp\'ecialisation co\"incide avec l'ordre de $L$. Les ferm\'es irr\'eductibles
sont les ensembles $\downarrow\! p$ pour $p$ \'el\'ement irr\'eductible de rencontre
(\emph{meet-irreducible}).
\end{exemple}

\begin{exercice}
Soit $P = \{a,b,c\}$ avec $a \leq b$, $a \leq c$, et $b,c$ incomparables. D\'ecrire
la topologie d'Alexandrov, les ferm\'es, les ferm\'es irr\'eductibles, et v\'erifier que
l'espace est $T_0$ mais ni $T_1$ ni sobre.
\end{exercice}

% ------------------------------------------------------------
\section{Topologie sup\'erieure d'un ordre}
% ------------------------------------------------------------

\begin{definition}[Topologie sup\'erieure]
Soit $(P, \leq)$ un ensemble ordonn\'e. La \emph{topologie sup\'erieure} (ou \emph{upper
topology}) est la topologie engendr\'ee par les compl\'ementaires des parties basses
principales~:
\[
  \tau_{\mathrm{up}} = \langle P \setminus \downarrow\! x : x \in P \rangle.
\]
C'est en g\'en\'eral plus grossi\`ere que la topologie d'Alexandrov.
\end{definition}

\begin{proposition}
La topologie sup\'erieure a le m\^eme ordre de sp\'ecialisation que la topologie
d'Alexandrov, \`a savoir l'ordre $\leq$ donn\'e. C'est la topologie la moins fine ayant
cet ordre de sp\'ecialisation.
\end{proposition}

\begin{proof}
Les ouverts de $\tau_{\mathrm{up}}$ sont des parties hautes, donc l'ordre de
sp\'ecialisation raffine $\leq$. R\'eciproquement, si $x \leq y$, alors tout
ouvert $P \setminus \downarrow\! z$ contenant $x$ (i.e., $x \not\leq z$) satisfait
$y \not\leq z$ (sinon $x \leq y \leq z$, contradiction), donc $y$ est dans cet ouvert.
Ainsi l'ordre de sp\'ecialisation est exactement $\leq$.
\end{proof}

% ------------------------------------------------------------
\section{Ensembles satur\'es}
% ------------------------------------------------------------

\begin{definition}[Ensemble satur\'e]
Un sous-ensemble $A$ d'un espace topologique est \emph{satur\'e} s'il est \'egal \`a
l'intersection de tous les ouverts le contenant~:
$A = \bigcap\{U \in \tau : A \subseteq U\}$.
\end{definition}

\begin{proposition}
Dans tout espace topologique~:
\begin{enumerate}
  \item Un sous-ensemble est satur\'e si et seulement s'il est une partie haute pour
    l'ordre de sp\'ecialisation.
  \item Tout ouvert est satur\'e.
  \item Le satur\'e d'un ensemble $A$ est $\uparrow\! A$.
  \item Un compact satur\'e est l'analogue correct d'un compact dans le cadre non-Hausdorff.
\end{enumerate}
\end{proposition}

\begin{proof}
(1) L'intersection de tous les ouverts contenant $A$ est l'ensemble des points $y$ tels
que tout ouvert contenant $A$ contient $y$. Or $y$ est dans tout ouvert contenant $x \in A$
si et seulement si $x \leq y$. Donc cette intersection est $\uparrow\! A$.
\end{proof}

% ------------------------------------------------------------
\section{Stably compact spaces (aper\c{c}u)}
% ------------------------------------------------------------

\begin{definition}[Espace stablement compact]
Un espace $T_0$ est \emph{stablement compact} s'il est~:
\begin{enumerate}
  \item compact,
  \item localement compact (tout point a une base de voisinages compacts satur\'es),
  \item sobre,
  \item \emph{coherent}~: l'intersection de deux compacts satur\'es est compacte satur\'ee.
\end{enumerate}
\end{definition}

\begin{remarque}
Les espaces stablement compacts jouent un r\^ole central dans la th\'eorie, car ils sont
exactement les espaces $T_0$ compacts dont le treillis des ouverts est un \emph{treillis
continu}. Nous y reviendrons au Chapitre~8.
\end{remarque}

% ------------------------------------------------------------
\section{Exercices}
% ------------------------------------------------------------

\begin{exercice}
Montrer qu'un espace est $T_1$ si et seulement si son ordre de sp\'ecialisation est
l'\'egalit\'e.
\end{exercice}

\begin{exercice}
Soit $X = \mathrm{Spec}(A)$ le spectre d'un anneau commutatif avec la topologie de Zariski.
D\'ecrire l'ordre de sp\'ecialisation. Montrer que le point g\'en\'erique de $X$ (s'il
existe) correspond \`a l'id\'eal premier minimal.
\end{exercice}

\begin{exercice}
Montrer que dans un espace $T_0$, l'adh\'erence d'un singleton $\overline{\{x\}}$ est
toujours un ferm\'e irr\'eductible.
\end{exercice}

\begin{exercice}
Construire un espace $T_0$ dans lequel il existe un ferm\'e irr\'eductible qui n'est
l'adh\'erence d'aucun singleton. (Indication~: consid\'erer un ensemble ordonn\'e sans
plus grand \'el\'ement, muni de la topologie sup\'erieure.)
\end{exercice}

\begin{exercice}
Soit $f\colon X \to Y$ une application continue ouverte entre espaces $T_0$.
Montrer que $f$ est croissante et que $f^{-1}(\uparrow\! y) = \uparrow\! f^{-1}(y)$
pour tout $y \in Y$.
\end{exercice}

\begin{exercice}[Topologie d'Alexandrov et limites directes]
Soit $(P_i, f_{ij})_{i \leq j}$ un syst\`eme direct de posets finis avec fonctions
croissantes. Montrer que la limite directe dans $\mathbf{Top}$ des espaces d'Alexandrov
correspondants est un espace dont l'ordre de sp\'ecialisation est la limite directe des
ordres.
\end{exercice}
